\documentclass[10pt,a4paper]{amsart}
\usepackage[margin=19mm]{geometry}
\usepackage{amsmath,amssymb,mathtools}
\usepackage[scr=rsfs,cal=euler]{mathalfa}
\usepackage{xfrac}
\usepackage[unicode,hidelinks,bookmarks=false]{hyperref}
\newcommand{\ie}{i.e.\ }
\newcommand{\cf}{cf.\ }
\newcommand{\resp}{respectively\ }
\newcommand{\Subsection}{\subsection}
\providecommand{\nobreakdash}{\nobreak\hbox{-}}
\setlength{\emergencystretch}{2em}
\multlinegap0pt
\usepackage{amssymb}
\newtheorem{theorem}{Theorem}
\title{A note on unitary invariance of Connes spectral distances of quantum states}
\author{Bing-Sheng Lin, Ji-Hong Wang, Zhi-Kang You}
\date{Source: arXiv:2605.13014v2 (CC BY 4.0)}
\begin{document}
\maketitle

\noindent\emph{Source and licence.} A note on unitary invariance of Connes spectral distances of quantum states. Version arXiv:2605.13014v2 (CC BY 4.0), distributed by arXiv under the Creative Commons Attribution 4.0 International licence, http://creativecommons.org/licenses/by/4.0/, as recorded in the arXiv licence metadata for this exact version and shown on the abstract page https://arxiv.org/abs/2605.13014v2. Full licence notice: \texttt{licenses/arxiv-2605.13014-CC-BY-4.0.txt}.\par\smallskip

\noindent\textit{Editorial scope.} Counted unit: the article's theorem theorem (source0.tex lines 1097--1110). The article's complete original proof of it is delivered with it (source0.tex lines 1112--1176, one \texttt{proof} environment of 65 lines). No mathematics is written, repaired or re-derived: the statement and the proof are the article's own lines, in the article's own notation, and no retained line is substituted or re-flowed.  Retained uncounted as the source-local context the proof needs: lines 1088--1091.  Nothing was omitted from any retained span, so the package carries no omission record.  No retained line contains a citation, so the wrapper carries no bibliography and no bibliography entry.  No retained line contains a figure environment, an image-inclusion command or drawing code, so no image is delivered and the package declares no asset dependency.  Editorial formatting only, in addition: an AMS \texttt{amsart} preamble that restates the article's own declarations and macros used by the retained lines, namely the environments \texttt{theorem} on separate counters, together with the standard packages those lines need.  The retained environments print in this document as Theorem 1 in order of appearance, whereas the article prints the same items with the numbering of its own full document.  The complete native source of the article as distributed in the arXiv e-print for this exact version is bundled unchanged as \texttt{sources/200471/source0.tex}.
  \begin{equation}\label{d4p}
    \mathcal{D}_4' = \frac{1}{4}  \sum_{i = 1}^3 \pm \sigma_i \otimes
    \sigma_{i'}
  \end{equation}

\begin{theorem}
  Consider the Dirac operator,
  \begin{equation}
    \mathcal{D}_{4^n} = \frac{1}{4}  \sum_{i = 1}^3 \sigma_{i_n, i_n'}
    \otimes \cdots \otimes \sigma_{i_1, i_1'},
  \end{equation}
  where $i_k, i_k'$ are any permutations of $\{1, 2, 3\}$, and $\sigma_{i_k,
  i_k'} = \pm \sigma_{i_k} \otimes \sigma_{i_k'}$. Then for any $2 \times 2$
  traceless Hermitian matrix $e \in \mathbb{M}_2 (\mathbb{C})$, we have
  \begin{equation}
    \|[\mathcal{D}_{4^n}, \pi_n (e)]\|_{op} = \|e\|_{op},
  \end{equation}
  where $\pi_n (e) =\mathbb{I}_{4^{n - 1}} \otimes \mathbb{I}_2 \otimes e$.
\end{theorem}

\begin{proof}
	When $n = 1$, the Dirac operator $\mathcal{D}_{4}$ is just $\mathcal{D}_4'$ (\ref{d4p}).

When $n = 2$, we have
\begin{equation}
	\mathcal{D}_{16} = \frac{1}{4}  \sum_{i = 1}^3 \sigma_{i_2, i_2'} \otimes
	\sigma_{i_1, i_1'},
\end{equation}
and
\begin{eqnarray}
	\|[\mathcal{D}_{16}, \pi_2 (e)]\|_{op} & = & \left\| \left[ \frac{1}{4} 
	\sum_{i = 1}^3 \sigma_{i_2, i_2'} \otimes \sigma_{i_1, i_1'}, \mathbb{I}_4
	\otimes \pi_1 (e) \right] \right\|_{op} \nonumber\\
	& = & \frac{1}{4} \left\| \sum_{i = 1}^3 \sigma_{i_2, i_2'} \otimes
	[\sigma_{i_1, i_1'}, \pi_1 (e)] \right\|_{op} . 
\end{eqnarray}
Using the unitary matrices $U_{\pm}$, one can permutate the Pauli matrices. So
one can only consider the following simplest case,
\begin{equation}
	\|[\mathcal{D}_{16}, \pi_2 (e)]\|_{op} = \frac{1}{4} \left\| \sum_{i = 1}^3
	\sigma_{i, i} \otimes [\sigma_{i, i}, \pi_1 (e)] \right\|_{op} .
\end{equation}
There is the following $4 \times 4$ unitary matrix
\begin{equation}
	U = \frac{1}{\sqrt{2}} \left( \begin{array}{cccc}
		1 & 0 & 0 & - 1\\
		0 & 1 & - 1 & 0\\
		0 & 1 & 1 & 0\\
		1 & 0 & 0 & 1
	\end{array} \right),
\end{equation}
and we have
\begin{eqnarray}
	\|[\mathcal{D}_{16}, \pi_2 (e)]\|_{op} 
	& = & \left\| (U \otimes \mathbb{I}_4) \left( \frac{1}{4}  \sum_{i = 1}^3
	\sigma_{i, i} \otimes [\sigma_{i, i}, \pi_1 (e)] \right) (U^{\dag} \otimes
	\mathbb{I}_4) \right\|_{op} \\
	& = & \left\| \left( \begin{array}{cccc}
		{}[D_1, \pi_1 (e)] & 0 & 0 & 0\\
		0 & [D_2, \pi_1 (e)] & 0 & 0\\
		0 & 0 & [D_3, \pi_1 (e)] & 0\\
		0 & 0 & 0 & [D_4, \pi_1 (e)]
	\end{array} \right) \right\|_{op}, \nonumber
\end{eqnarray}
where
\begin{equation}
	D_k = \frac{1}{4}  \sum_{i = 1}^3 \sigma_{i, i} = \frac{1}{4}  \sum_{i =
		1}^3 \pm \sigma_i \otimes \sigma_i, \quad k = 1, 2, 3, 4,
\end{equation}
and different $D_k$ has different signs of some terms, but they all have the
same relation
\begin{equation}
	\|[D_k, \pi_1 (e)]\|_{op} =\|[\mathcal{D}_4, \pi_1 (e)]\|_{op}= \|e\|_{op}, \quad k = 1, 2, 3, 4 .
\end{equation}
So we have
\begin{equation}
	\|[\mathcal{D}_{16}, \pi_2 (e)]\|_{op} = \|[\mathcal{D}_4, \pi_1 (e)]\|_{op}
	= \|e\|_{op} . 
\end{equation}
Using the similar methods, one can easily find that, for any $n \geqslant 1$, there is
\begin{equation}
	\|[\mathcal{D}_{4^n}, \pi_n (e)]\|_{op} = \|[\mathcal{D}_4, \pi_1 (e)]\|_{op}
	= \|e\|_{op} .
\end{equation}
\end{proof}

\end{document}
